\section{Consolidation Cycles}\label{app:consolidation}

The Hopfield store $\mathcal{S}$ holds a set of normalised embedding vectors. Two consolidation cycles run on user request through the MCP tools \texttt{pce.store.consolidate\_sws} and \texttt{pce.store.consolidate\_rem}; both are deterministic given a seed.

\subsection{Slow-wave-sleep abstraction}

Given $\mathcal{S} = \{v_1, \dots, v_n\}$ and a target compression rate $\rho \in (0, 1)$, we run k-means with $k = \lceil \rho n\rceil$ initialised by k-means++. Cosine distance is used; the centroid is the unit-normalised mean. The centroids \emph{replace} the cluster members in $\mathcal{S}$: this is the abstraction step. Numerical safeguard: cosine distances are clipped to $[0, 2]$ before normalisation in the k-means++ probability update (we hit a \texttt{Probabilities are not non-negative} crash here in Phase~8 before adding the clamp; see commit \texttt{a27e886}).

\subsection{REM-style replay}

Given $\mathcal{S}$ and a temperature $T$, sample $m$ items by a Metropolis chain on $\mathcal{S}$ with energy $E(v) = -\log p(v)$ where $p$ is the empirical density estimated from cluster sizes. Each accepted item is perturbed by $v\!+\!\epsilon$, $\epsilon\sim\mathcal{N}(0, \sigma I)$, and re-stored. This raises the salience of rare-but-coherent attractors without erasing dense modes. Defaults: $T{=}0.7$, $\sigma{=}0.05$, $m=\min(n/4, 32)$.

\subsection{Why both}

SWS abstracts; REM diversifies. A single cycle of either alone collapses the store toward the data mode (SWS) or noise (REM). In our cascade we run them in alternation: SWS after a creative session, REM before the next session. The order is recorded in the audit log so a downstream user can replay both cycles deterministically.
